\documentclass[12pt]{scrartcl}

\usepackage[utf8]{inputenc}
\usepackage[T1]{fontenc}
\usepackage[ngerman]{babel}
\usepackage{lmodern}
\usepackage{graphicx}
\usepackage[pdftex,hyperref,dvipsnames]{xcolor}
\usepackage{listings}
\usepackage[a4paper,lmargin={2cm},rmargin={2cm},tmargin={3.5cm},bmargin = {2.5cm},headheight = {4cm}]{geometry}
\usepackage{amsmath,amssymb,amstext,amsthm}
% \usepackage[lined,algonl,boxed]{algorithm2e} 
\usepackage{algpseudocode}
\usepackage{tikz}
\usepackage{pgfplots}
\usepackage{hyperref}
\usepackage{url}
\usepackage{tcolorbox}
\usepackage[inline]{enumitem}
\usepackage[headsepline]{scrlayer-scrpage} 

\usepackage[scale=0.55, lineVdist=1em]{secstack}
\pagestyle{scrheadings} 
\usetikzlibrary{automata,positioning}

\usepgfplotslibrary{external}
%\tikzexternalize
\pgfplotsset{width=10cm,compat=1.9}
\pgfplotsset{every axis/.style={
    axis line style = thick,
    grid = both,
    tick pos = left,
    tick align = outside,
    tick style = {thick,black},
    fill = blue,
	line width = 2pt,
}}

\lstset{
  basicstyle=\ttfamily\small,
  numbers=left,
  numberstyle=\tiny,
  stepnumber=1,
  numbersep=8pt,
  frame=single,
  breaklines=true,
  columns=fullflexible
}

\lstset{
  language=C,
  basicstyle=\ttfamily\small,
  keywordstyle=\color{blue},
  commentstyle=\color{gray},
  stringstyle=\color{green!60!black},
  numbers=left,
  numberstyle=\tiny,
  stepnumber=1,
  numbersep=8pt,
  frame=single,
  breaklines=true,
  columns=fullflexible
}

\renewcommand{\theenumi}{(\alph{enumi})}
\renewcommand{\labelenumi}{\text{\theenumi}}
\renewcommand{\labelenumii}{(\roman{enumii})} 

\ohead{Marvin\\
       Jan Theil (3794698)}

\ihead{System and Web Security}

\newcounter{exnum}
\newcommand{\exercise}[1]{\section*{Problem \theexnum\stepcounter{exnum}: #1}}

\begin{document}

\setcounter{exnum}{1}

\exercise{Disassembling with gdb}

The first letters of each team member's given name are $M$ and $J$.
This results in the following code:

\begin{lstlisting}
int main() {
  int a;
  int b;
  int c;
  a = 'M' + 'J';
  b = 3;
  c = a+b;
  return 0;
}
\end{lstlisting}

\texttt{gdb command disassemble main} outputs:
\begin{lstlisting}[language={}]
Dump of assembler code for function main:
0x08048344 <main+0>:	lea    0x4(%esp),%ecx
0x08048348 <main+4>:	and    $0xfffffff0,%esp
0x0804834b <main+7>:	pushl  -0x4(%ecx)
0x0804834e <main+10>:	push   %ebp
0x0804834f <main+11>:	mov    %esp,%ebp
0x08048351 <main+13>:	push   %ecx
0x08048352 <main+14>:	sub    $0x10,%esp
0x08048355 <main+17>:	movl   $0x97,-0x10(%ebp)
0x0804835c <main+24>:	movl   $0x3,-0xc(%ebp)
0x08048363 <main+31>:	mov    -0xc(%ebp),%eax
0x08048366 <main+34>:	add    -0x10(%ebp),%eax
0x08048369 <main+37>:	mov    %eax,-0x8(%ebp)
0x0804836c <main+40>:	mov    $0x0,%eax
0x08048371 <main+45>:	add    $0x10,%esp
0x08048374 <main+48>:	pop    %ecx
0x08048375 <main+49>:	pop    %ebp
0x08048376 <main+50>:	lea    -0x4(%ecx),%esp
0x08048379 <main+53>:	ret    
End of assembler dump.
\end{lstlisting}

\begin{itemize}
    \item Line 5 (\texttt{a = 'M' + 'J'}) corresponds to \texttt{0x08048355, line 9}:\\
          \texttt{'M' + 'J' = 0x4D + 0x4A = 0x97} and in this line \texttt{0x97} is written to 
          \texttt{ebp - 0x10} which seems to be the location of local variable \texttt{a}.
    \item Line 6 (\texttt{b = 3}) corresponds to \texttt{0x0804835c, line 10}:\\
          Same reasoning as above. \texttt{3 = 0x3} is written to \texttt{ebp - 0xC}
          which seems to be the location of local variable \texttt{b}.
    \item Line 7 (\texttt{c = a+b}) corresponds to \texttt{0x08048363 - 0x08048369, line 11 - 13}:\\
          First the value of \texttt{ebp - 0xC = variable b} is written to register \texttt{eax}. Then
          the value of \texttt{ebp - 0x10 = variable a} is added to \texttt{eax} resulting in \texttt{eax = a + b}.
          Finally, the result of the addition is written to texttt{ebp - 0x8} which seems to be the location of local variable \texttt{c}.
\end{itemize}


\exercise{x86 Assembly Code (I)}

% #include <stdio.h>

% void myfunction(int a) {
%         long some_long;
%         char * some_strings[2];
%         some_long = 42;
%         some_strings[0] = "foo";
%         some_strings[1] = "bar";
%         some_long -= 19;
%         printf("%ld\n%d\n%s\n", some_long, a, some_strings[0]);
% }
% int main () {
%         int x = 0;
%         myfunction(x);
% }


\begin{lstlisting}[language={}, caption=Assembler code for main]
Dump of assembler code for function main:
0x080483b6 <main+0>:	lea    0x4(%esp),%ecx  # copy address ESP - 0x4 to ECX
0x080483ba <main+4>:	and    $0xfffffff0,%esp  # alignment
0x080483bd <main+7>:	pushl  -0x4(%ecx)  # push ECX - 0x4 to the stack
0x080483c0 <main+10>:	push   %ebp  # push EBP to the stack
0x080483c1 <main+11>:	mov    %esp,%ebp  # EBP = ESP
0x080483c3 <main+13>:	push   %ecx  # push ECX to the stack
0x080483c4 <main+14>:	sub    $0x14,%esp  # subtract 0x14 from ESP
0x080483c7 <main+17>:	movl   $0x0,-0x8(%ebp)  # write 0x0 to EBP - 0x8
0x080483ce <main+24>:	mov    -0x8(%ebp),%eax  # write EBP - 0x8 to EAX 
0x080483d1 <main+27>:	mov    %eax,(%esp)  # write EAX to ESP
0x080483d4 <main+30>:	call   0x8048374 <myfunction>  # call myfunction: push 0x8048374 to stack and jump
0x080483d9 <main+35>:	add    $0x14,%esp  # add 0x14 to ESP
0x080483dc <main+38>:	pop    %ecx  # pop ECX from stack
0x080483dd <main+39>:	pop    %ebp  # pop abp from stack
0x080483de <main+40>:	lea    -0x4(%ecx),%esp  # copy address ECX - 0x4 to ESP
0x080483e1 <main+43>:	ret  # return: pop instruction pointer from stack
End of assembler dump.
\end{lstlisting}

\begin{lstlisting}[language={}, caption=Assembler code for myfunction]
Dump of assembler code for function myfunction:
0x08048374 <myfunction+0>:	push   %ebp  # Push EBP to stack
0x08048375 <myfunction+1>:	mov    %esp,%ebp  # ESP = EBP: stack pointer is new fram pointer
0x08048377 <myfunction+3>:	sub    $0x28,%esp  # ESP = ESP - 0x28 
0x0804837a <myfunction+6>:	movl   $0x2a,-0x4(%ebp)  # EBP - 0x4 = 0x2a: Init some_long with 42
0x08048381 <myfunction+13>:	movl   $0x80484b0,-0xc(%ebp)  # EBP - 0xC = 0x80484b0: some_strings[...] = "..."
0x08048388 <myfunction+20>:	movl   $0x80484b4,-0x8(%ebp)  # EBP - 0x8 = 0x80484b4: some_strings[...] = "..."
0x0804838f <myfunction+27>:	subl   $0x13,-0x4(%ebp)  # EBP - 0x4 = (EBP - 0x4) - 0x13: some_long -= 19
0x08048393 <myfunction+31>:	mov    -0xc(%ebp),%eax  # EAX = EBP - 0xC: Save some_string[...] to EAX
0x08048396 <myfunction+34>:	mov    %eax,0xc(%esp)  # ESP + 0xC = EAX: Write EAX (=some_string[...]) to the stack
0x0804839a <myfunction+38>:	mov    0x8(%ebp),%eax  # EAX = EBP + 0x8: Save some_string[...] to EAX
0x0804839d <myfunction+41>:	mov    %eax,0x8(%esp)  # ESP + 0x8 = EAX: Write EAX (=some_string[...]) to the stack
0x080483a1 <myfunction+45>:	mov    -0x4(%ebp),%eax  # EBP - 0x4 = EAX: Save some_long to EAX
0x080483a4 <myfunction+48>:	mov    %eax,0x4(%esp)  # ESP + 0x4 = EAX: Write EAX (=some_long) to the stack
0x080483a8 <myfunction+52>:	movl   $0x80484b8,(%esp)  # ESP = 0x80484b8: printf string pointer
0x080483af <myfunction+59>:	call   0x80482d8 <printf@plt>  # print call
0x080483b4 <myfunction+64>:	leave  # restore frame: ESP = EBP, pop EBP
0x080483b5 <myfunction+65>:	ret  # retrun
End of assembler dump.
\end{lstlisting}

\exercise{x86 Assembly Code (II)}

\begin{lstlisting}[language={}]
Dump of assembler code for function blub:
# Frame pointer setup
0x080483c4 <blub+0>: push %ebp
0x080483c5 <blub+1>: mov %esp,%ebp
0x080483c7 <blub+3>: sub $0x18,%esp  # Allocates 0x18 bytes on the stack

# Actual function
0x080483ca <blub+6>: movl $0x17,-0x4(%ebp)  # Init EBP - 0x4 (some local variable, let's say "a") with 0x17
0x080483d1 <blub+13>: mov -0x4(%ebp),%eax  # EAX = EBP - 0x4 = a
0x080483d4 <blub+16>: sub $0xf,%eax  # EAX = EAX - 0x: a - 0xf
0x080483d7 <blub+19>: mov %eax,-0x8(%ebp)  # EBP - 0x8 = EAX: save EAX to some other local variable ("b")
0x080483da <blub+22>: mov -0x8(%ebp),%eax  # EAX = EBP - 0x8 = b
0x080483dd <blub+25>: mov %eax,0x4(%esp)  # ESP - 0x4 = EAX: Write value of EAX = b to stack
0x080483e1 <blub+29>: movl $0x80484d0,(%esp)  # Write some printf string to stack
0x080483e8 <blub+36>: call 0x80482f8 <printf@plt>  # call printf

# Return
0x080483ed <blub+41>: leave
0x080483ee <blub+42>: ret
End of assembler dump.
\end{lstlisting}

The following c function is functionally equivalent (+ added main()):
\begin{lstlisting}[language={C}]
#include <stdio.h>

void blob() {
    int a;
    int b;
    a = 0x17;
    b = a - 0xf;
    printf("b = %x", b);
}

int main() {
    blob();
    return 0;
}
\end{lstlisting}

% Dump of assembler code for function blob:
% 0x08048374 <blob+0>:	push   %ebp
% 0x08048375 <blob+1>:	mov    %esp,%ebp
% 0x08048377 <blob+3>:	sub    $0x18,%esp
% 0x0804837a <blob+6>:	movl   $0x17,-0x8(%ebp)
% 0x08048381 <blob+13>:	mov    -0x8(%ebp),%eax
% 0x08048384 <blob+16>:	sub    $0xf,%eax
% 0x08048387 <blob+19>:	mov    %eax,-0x4(%ebp)
% 0x0804838a <blob+22>:	mov    -0x4(%ebp),%eax
% 0x0804838d <blob+25>:	mov    %eax,0x4(%esp)
% 0x08048391 <blob+29>:	movl   $0x8048490,(%esp)
% 0x08048398 <blob+36>:	call   0x80482d8 <printf@plt>
% 0x0804839d <blob+41>:	leave  
% 0x0804839e <blob+42>:	ret    
% End of assembler dump.


\exercise{Two’s Complement}
\begin{align*}
x &= 'M' + 'J' = 0x4D + 0x4A = \\
  &= 0x97 = 151\\
\implies s &= (x \bmod 127) + 128 \\
           &= (151 \bmod 127) + 128 \\
           &= (24) + 128 = 152
\end{align*}

Binary representation of $s$ as an unsigned (8-bit) number: $s = 152 = 10011000_2$

Interpretation of $s = 10011000_2$ as a signed (8-bit) number:
\begin{itemize}
    \item MSB $\mathbf{1}0011000_2$ is $1$ which inidcates a negative value
    \item Calculate 2's complement: $10011000_2 \overset{\text{invert}}{\to} 01100111_2 \overset{+1}{\to} 01101000_2$
    \item $01101000_2 = 104 \implies s = -104$
\end{itemize} 

\end{document}